\section{Band Power \& the Freeze Index}
\label{sec:freeze-index}

The previous stage handed us a Welch power spectral density (PSD): a curve that
says \emph{how much} of the leg's motion lives at each frequency. A spectrum is rich,
but a state machine cannot act on a whole curve --- it needs a single number it can
compare against a threshold. This section reduces the spectrum to exactly that number,
the \textbf{Freeze Index}, by first collapsing the curve into a handful of
\emph{band powers} and then taking a ratio of two of them. The ratio is the small,
quietly clever idea that lets a \$30 board tell a freeze from ordinary walking.

\begin{intuition}
Walking and freezing both shake the ankle --- the difference is \emph{where} in the
frequency spectrum the energy sits. A walking stride is a slow, heavy rhythm
(roughly one to three steps a second); a freeze is a fast, fluttering tremor of the
feet \emph{in place} at 3--8\,Hz. So we do not ask ``how hard is the leg moving?''
(walking wins that contest) --- we ask ``what \emph{fraction} of the shaking is fast?''.
The Freeze Index is precisely that fraction expressed as a ratio, and it flips
decisively the moment a freeze begins.
\end{intuition}

\subsection{From a spectrum to a band power}

The PSD $S(f)$ tells us the motion's power density (power per hertz) at every
frequency bin. We rarely care about one bin in isolation; we care about a whole
\emph{band} of physiologically meaningful frequencies. The \textbf{band power} over a
band $[f_1,f_2]$ is simply the PSD summed across the bins that fall inside it:
%
\begin{equation}
P[f_1,f_2] \;=\; \sum_{k\,:\, f_1 \le f_k \le f_2} S(f_k)\,,
\qquad f_k = k\,\Delta f,\quad \Delta f = \frac{f_s}{N} = \frac{64}{256}=\SI{0.25}{\hertz}.
\label{eq:bandpower}
\end{equation}
%
Because the window is $N=256$ samples (\SI{4.0}{\second}) at $f_s = \SI{64}{\hertz}$,
the bins are spaced $\Delta f = \SI{0.25}{\hertz}$ apart --- fine enough to slice the
spectrum cleanly into the three bands that matter clinically.

\begin{table}[H]
\centering
\caption{The three frequency bands Cadence reads off the mean-removed magnitude
spectrum. Each band is a contiguous run of \SI{0.25}{\hertz} PSD bins summed via
Eq.~\eqref{eq:bandpower}.}
\label{tab:bands}
\small
\begin{tabular}{@{}llp{7.6cm}@{}}
\toprule
\textbf{Band} & \textbf{Range} & \textbf{What lives there} \\
\midrule
Locomotor & \SI{0.5}{}--\SI{3}{\hertz} & the walking rhythm: stride and step cadence,
the heavy slow swing of normal gait. \\
Freeze & \SI{3}{}--\SI{8}{\hertz} & the tell-tale rapid trembling/shuffling of the
feet \emph{in place} during a freeze episode. \\
Rest-tremor & \SI{4}{}--\SI{6}{\hertz} & Parkinsonian rest tremor; reported on the
dashboard as context, not used in the index. \\
\bottomrule
\end{tabular}
\end{table}

\begin{keyidea}
Every band power in Eq.~\eqref{eq:bandpower} is taken on the \emph{mean-removed
magnitude} $a=\norm{(a_x,a_y,a_z)}$, with its average subtracted. Combining the three
axes makes the measure orientation-invariant (the sleeve can rotate on the ankle);
subtracting the mean strips out gravity's $\sim\!1\,g$ DC offset so it never
contaminates the low band. We are left with pure \emph{motion} power, ready to slice.
\end{keyidea}

\subsection{The Freeze Index: a ratio of two bands}

Moore \emph{et al.} (2008) defined the Freeze Index as the power in the freeze band
divided by the power in the locomotor band:
%
\begin{equation}
\FI \;=\; \frac{P[\,\SI{3}{}\text{--}\SI{8}{\hertz}\,]}
              {P[\,\SI{0.5}{}\text{--}\SI{3}{\hertz}\,] + \varepsilon}\,,
\qquad \varepsilon = 10^{-9}.
\label{eq:freezeindex}
\end{equation}
%
The numerator is the energy of the fast freeze tremor; the denominator is the energy
of the slow walking rhythm. A tiny guard $\varepsilon = 10^{-9}$ is added to the
denominator so a perfectly still leg --- whose locomotor power can round to zero ---
never triggers a divide-by-zero. It is far below any real band power, so it changes no
genuine reading; it only tames the $0/0$ corner case.

\begin{mathsbox}
Read Eq.~\eqref{eq:freezeindex} as a competition between two bands.
\begin{itemize}
\item \textbf{Walking:} the slow stride dominates, so the denominator
$P[0.5\text{--}3]$ is large and the numerator is small $\Rightarrow \FI < 1$ to
$\sim\!1$.
\item \textbf{Freezing:} the fast 3--8\,Hz tremor dominates, so the numerator swells
and the denominator collapses $\Rightarrow \FI \gg 1$.
\end{itemize}
The index does not measure \emph{how much} the leg moves --- it measures the
\textbf{shape} of the spectrum, the balance of fast energy against slow energy. That
is why a ratio, not a difference, is the right tool: a difference would still be
dominated by amplitude.
\end{mathsbox}

\subsubsection*{Why a ratio? Scale-invariance.}

Suppose we measured acceleration in milli-$g$ on the host but in $\si{\meter\per\square\second}$
on-device --- a constant factor $c$ between the two unit systems. Power scales with the
\emph{square} of amplitude, so \emph{every} PSD bin, and hence both band powers, pick
up the same factor $c^2$:
%
\begin{equation}
\FI \;=\; \frac{c^2\,P[\,\SI{3}{}\text{--}\SI{8}{}\,]}
                {c^2\,P[\,\SI{0.5}{}\text{--}\SI{3}{}\,]}
       \;=\; \frac{P[\,\SI{3}{}\text{--}\SI{8}{}\,]}{P[\,\SI{0.5}{}\text{--}\SI{3}{}\,]}\,.
\label{eq:scaleinv}
\end{equation}
%
The factor cancels top and bottom. The Freeze Index is therefore
\textbf{scale-invariant}: it returns the same value whether the firmware works in
milli-$g$ or $\si{\meter\per\square\second}$, and the same value for a heavy stamp or a
light shuffle of the \emph{same spectral shape}. Amplitude is normalised away by
construction, leaving only the quantity we actually care about --- the proportion of
fast tremor.

\begin{figure}[H]
\centering
\resizebox{0.95\textwidth}{!}{%
\begin{tikzpicture}
  \begin{groupplot}[
      group style={group size=2 by 1, horizontal sep=2.2cm},
      width=8.2cm, height=6.0cm,
      xmin=0, xmax=9, ymin=0, ymax=1.05,
      xlabel={frequency (Hz)}, ylabel={PSD $S(f)$ (normalised)},
      xtick={0,2,4,6,8}, ytick=\empty,
      tick label style={font=\small}, label style={font=\small},
      title style={font=\bfseries\small, color=ink},
      axis line style={mute}, every axis/.append style={font=\small},
    ]
    % ---------- WALKING ----------
    \nextgroupplot[title={Walking: slow band tall}]
      % locomotor band shade 0.5-3
      \fill[grn!18] (axis cs:0.5,0) rectangle (axis cs:3,1.05);
      % freeze band shade 3-8
      \fill[freeze!12] (axis cs:3,0) rectangle (axis cs:8,1.05);
      % spectrum: big peak near 1.6 Hz, small ripple in freeze band
      \addplot[navy, very thick, smooth, domain=0.2:9, samples=180]
        { exp(-((x-1.6)^2)/0.55) + 0.07*exp(-((x-4.6)^2)/1.2) };
      \node[grn!60!black, font=\scriptsize\bfseries] at (axis cs:1.7,0.93) {loco};
      \node[freeze!70!black, font=\scriptsize\bfseries] at (axis cs:5.5,0.18) {freeze};
      \node[anchor=south east, font=\scriptsize\bfseries, navy]
        at (axis cs:8.8,0.40) {$\FI \approx 0.5\text{--}1$};
    % ---------- FREEZING ----------
    \nextgroupplot[title={Freezing: fast band tall}]
      \fill[grn!18] (axis cs:0.5,0) rectangle (axis cs:3,1.05);
      \fill[freeze!12] (axis cs:3,0) rectangle (axis cs:8,1.05);
      % spectrum: small loco, big peak near 5 Hz
      \addplot[navy, very thick, smooth, domain=0.2:9, samples=180]
        { 0.16*exp(-((x-1.6)^2)/0.55) + exp(-((x-5.0)^2)/0.9) };
      \node[grn!60!black, font=\scriptsize\bfseries] at (axis cs:1.7,0.30) {loco};
      \node[freeze!70!black, font=\scriptsize\bfseries] at (axis cs:5.0,0.93) {freeze};
      \node[anchor=south east, font=\scriptsize\bfseries, freeze]
        at (axis cs:8.8,0.55) {$\FI \gg 1$};
  \end{groupplot}
\end{tikzpicture}}
\caption{The Freeze Index reads the \emph{shape} of the spectrum. In walking (left)
the energy piles up in the green locomotor band (0.5--3\,Hz), so the ratio
freeze$/$loco stays near or below one. In a freeze (right) the energy migrates into
the red 3--8\,Hz band, the denominator collapses and the numerator swells, and the
ratio flips far above one. Amplitude (overall height) is irrelevant --- only the
balance between the two shaded areas drives $\FI$.}
\label{fig:two-spectra}
\end{figure}

\subsection{What it does on real data}

Figure~\ref{fig:two-spectra} is the cartoon; Figure~\ref{fig:fi-real} is the truth.
The capture \code{capture1.csv} (\SI{64}{\hertz}, \SI{41.4}{\second}) walks the ankle
through four scripted phases --- still, walk, freeze, walk --- and we compute $\FI$ on
every \SI{4}{\second} window. The contrast is exactly what
Eq.~\eqref{eq:freezeindex} predicts.

\begin{table}[H]
\centering
\caption{Real per-phase Freeze Index on \code{capture1.csv}. The freeze phase is an
order of magnitude above quiet standing and ordinary walking, and its peak windows
soar past 20.}
\label{tab:fi-phases}
\small
\begin{tabular}{@{}lccc@{}}
\toprule
\textbf{Phase} & \textbf{Mean $\FI$} & \textbf{Median $\FI$} & \textbf{Max $\FI$} \\
\midrule
Still (rest) & $\sim 1.1$ & --- & $1.71$ \\
Walk         & $\sim 1.2$ & --- & $2.08$ \\
\textbf{Freeze} & $\mathbf{6.89}$ & $4.45$ & $\mathbf{21.35}$ \\
\bottomrule
\end{tabular}
\end{table}

Two features deserve a careful eye. First, still and walk sit close together
($\sim\!1.1$ versus $\sim\!1.2$): \textbf{walking is not what trips the index}. Both
are quiet in the freeze band, so both keep $\FI$ near one. Second, the gap to the
freeze phase is enormous --- a mean of $6.89$ and a maximum of $21.35$, roughly six
times higher than either calm phase. That separation is what makes a single scalar
threshold viable.

\begin{figure}[H]
\centering
\includegraphics[width=0.82\linewidth]{fi_per_phase}
\caption{Measured Freeze Index per scripted phase of \code{capture1.csv} (still / walk
/ freeze / walk). Bars are the mean $\FI$ in each phase; the grey ticks mark the
per-phase maximum (the freeze phase peaks at $21.35$, off the top of the calm bars).
Still and walk hug the deployed threshold line; the freeze phase towers over it. The
gap, not the absolute height, is what a detector exploits.}
\label{fig:fi-real}
\end{figure}

\begin{intuition}
Notice how flat the two calm bars are next to the freeze bar. If we had used a
\emph{difference} of band powers instead of a ratio, a vigorous walk could rack up
large freeze-band energy in absolute terms and look alarming. The ratio refuses to be
fooled: a vigorous walk \emph{also} has large locomotor energy, so the fraction stays
low. Only when the spectral balance genuinely tips toward 3--8\,Hz --- a real freeze
--- does the number climb.
\end{intuition}

\subsection{A scalar the state machine can act on}

The Freeze Index turns a four-second spectrum into one comparable number, refreshed
every two seconds as the window hops. Downstream, a window is flagged as a candidate
freeze when $\FI$ crosses a trip line (the deployed value, $1.815$, sits just above
the calm-phase maxima of $1.71$ and $2.08$ so that genuine freezes clear it). On the
live dashboard the same quantity drives the \emph{Freeze Index} readout (walk
$\sim\!0.5$--$1$; freeze $>2$, up to $20{+}$) and the \emph{Freeze Margin} meter, which
expresses $\FI$ as a percentage of that $1.815$ trip line. One ratio, computed once
per window, feeds the verdict, the margin meter, and the threshold sweep alike.

It is worth being explicit about what the index does \emph{not} settle on its own. A
large $\FI$ means the spectrum is freeze-shaped, but a near-motionless leg can produce
a freeze-shaped spectrum out of pure sensor noise --- a few micro-$g$ of jitter, mostly
high-frequency, divided by an almost-empty locomotor band. The $\varepsilon$ guard
keeps that quotient finite, but it does not keep it \emph{small}. The Freeze Index
alone would therefore raise false alarms whenever the wearer simply stands still. For
that reason Cadence never trusts $\FI$ in isolation: it is paired with a separate
\emph{movement-energy} gate that asks whether the leg is moving at all before the ratio
is allowed to vote. That gate, the threshold sweep that sets the $1.815$ trip line, and
the two-window debounce are the subject of the sections that follow; here it is enough
that the spectrum has been distilled into a single, honest, scale-invariant number.

\begin{takeaway}
The Freeze Index is one scale-invariant ratio --- fast-band power over slow-band power
--- that ignores how \emph{hard} the leg moves and reports only the \emph{shape} of its
motion. On real capture data it stays near $1.1$--$1.2$ for standing and walking but
leaps to a mean of $6.89$ (max $21.35$) during a freeze, giving the detector a clean,
single-number signal to act on.
\end{takeaway}
